\begin{statementbox}
	\textbf{\textcolor{CStatement}{先想一个问题}}
	已知 $2^3=8$。如果只知道 $2^{\,?}=8$，怎样找出这个“？”？
	对数就是帮我们找指数的：$\log_2 8=3$。

\begin{center}
\begin{tikzpicture}[font=\small,>=stealth]
\draw[draw=teal!50!black,fill=green!3,rounded corners=3pt] (0,0) rectangle (13,6);
\node at (6.5,5.45) {同一条指数曲线：已知高度找指数，已知指数找高度};
\begin{scope}[shift={(1.3,1.2)},x=1.1cm,y=.4cm]
\draw[gray,->] (-.6,0)--(3.8,0) node[right] {$x$};\draw[gray,->] (0,0)--(0,9.5) node[above] {$y$};
\draw[teal!70!black,very thick,domain=-.5:3.2,samples=70] plot(\x,{2^\x});
\draw[->,blue!70!black,thick] (0,8)--(2.95,8);
\draw[->,orange!85!black,thick] (3,7.9)--(3,.15);
\fill[teal!70!black] (3,8) circle(2.7pt);
\node[above left] at (3,8.3) {$(3,8)$};\node[left,blue!70!black] at (0,8) {$N=8$};
\node[below,orange!85!black] at (3,-.05) {$t=3$};\node[left] at (0,1) {$1$};
\end{scope}
\node[anchor=west,align=left,text width=5.4cm] at (6.8,3.1) {$y=2^x$ 上只有一个点的高度是 $8$。\\[10pt]横着读：$\log_2 8=3$。\\[10pt]竖着读：$2^3=8$。\\[10pt]两次读取的是同一个点。};
\node at (6.5,.45) {先找指数再算高度，或先算高度再找指数，都回到原数。};
\end{tikzpicture}

\end{center}

	\medskip
	\textbf{\textcolor{CStatement}{使用条件}}
	\begin{description}[style=nextline, leftmargin=2.8em, labelsep=0.5em, itemsep=0.35em]
		\item[\textbf{底数：}] $a>0$，且 $a\neq 1$。
		\item[\textbf{真数：}] $N>0$（只针对下面第一式中的 $N$）。
	\end{description}

	\textbf{\textcolor{CStatement}{两个核心结论}}
	\begin{description}[style=nextline, leftmargin=2.8em, labelsep=0.5em, itemsep=0.45em]
		\item[\textbf{先取对数，再算指数：}]
			\[
				\boxed{a^{\log_a N}=N}\qquad (N>0)
			\]
			例如 $2^{\log_2 8}=2^3=8$。先找到“几次方”，再算回原来的数。
		\item[\textbf{先算指数，再取对数：}]
			\[
				\boxed{\log_a(a^x)=x}\qquad (x\in\mathbb{R})
			\]
			例如 $\log_2(2^3)=\log_2 8=3$。已知算了几次方，再问指数，就得到 $3$。
	\end{description}

	\textbf{\textcolor{CStatement}{一眼记住}}
	\textbf{同底才能直接还原；对数里边的数必须大于 $0$。}
\end{statementbox}
